\section{Introduction}
\label{sec:introduction}

Euclid's theorem --- proved in Euclid's \emph{Elements}, Book IX,
Proposition 20 \cite{euclidelements} --- states that there are infinitely
many prime numbers. Euclid's original presentation uses a common multiple
of the assigned primes plus one. This article formalizes the product
(primorial) specialization of that construction:

\begin{quote}
Given any finite list of primes $p_1, p_2, \dots, p_k$, let
$N = p_1 \cdot p_2 \cdot \dots \cdot p_k + 1$.
Then $N$ is either prime itself, or has a prime divisor $d$ that is not
among $p_1, \dots, p_k$.
In either case, a new prime is found, proving the list cannot contain all
primes.
\end{quote}

In this article, we formalize and verify this proof using
\href{https://epfl-lara.github.io/stainless/intro.html}{Scala Stainless}
\cite{hamza2019systemfr}, a verification framework for pure Scala programs.
Our approach follows the zero-prior-knowledge methodology established in
earlier articles: modular arithmetic \cite{mata2026modulo}, lists
\cite{mata2026lists}, and prime utilities are all defined from scratch and
verified independently.

This article verifies:

\begin{itemize}
  \item Primorial-plus-one coprime to all list primes (Section~3.1).
  \item New prime found via the Euclid construction (Section~3.2).
  \item The new prime is not in the original list (Section~3.3).
  \item Euclid's theorem: primes are infinite (Section~3.4).
  \item Supporting verified prime lemmas (Section~4).
\end{itemize}
